% ============================================================================
% Section IV: Natural-Gas Network
% ============================================================================

\section{Natural-Gas Network}
\label{sec:gas}

Gas-fired units connect the power system to a pipeline network whose
pressures, compressor stations and line pack limit how much fuel can
reach each plant and when.  \gtopt{} represents this network with six
native elements (gas nodes, supplies, demands, pipes, compressors and
gas-turbine links) that are assembled into the same LP as the power
system of Section~\ref{sec:formulation}.  The model is steady state per
block, linear in the squared pressure, and contains no binary variable,
so it can be used by every solution method of
Section~\ref{sec:methods}.  Its compressor model follows
Saldarriaga~\cite{saldarriaga2018phd}, and its pipe model the
squared-pressure Weymouth form of De~Wolf and
Smeers~\cite{dewolf2000gas} and Menon~\cite{menon2005gas}.

\subsection{Sets, Units and Variables}

\begin{itemize}
  \item $\setN^{\mathrm{g}}$: gas nodes, indexed by $i, j$;
  \item $\mathcal{A}$: pipes, indexed by $a$, from inlet $i_a$ to
        outlet $j_a$;
  \item $\mathcal{C}$: compressor stations, indexed by $c$, from
        suction $i_c$ to discharge $j_c$;
  \item $\mathcal{K}$: gas supplies (fields, pipeline imports, LNG
        regasification), indexed by $k$, at node $i_k$;
  \item $\setD^{\mathrm{g}}$: gas demands, indexed by $d$, at node $i_d$;
  \item $\setG^{\mathrm{g}} \subseteq \setG$: gas-fired generators,
        each linked to a delivery node $i_g$.
\end{itemize}
Flows are in thousands of standard cubic metres per hour
[km\textsuperscript{3}/h], pressures in bar (absolute), line pack in
km\textsuperscript{3}, gas costs in \$/km\textsuperscript{3} and heat
rates in km\textsuperscript{3}/MWh.  The node pressure $\varrho_i$
enters the LP only through its square,
\begin{equation}
\label{eq:gas_pi}
\varpi_{i,s,t,b} = \varrho_{i,s,t,b}^2
  \in \bigl[\underline{\varrho}_i^{\,2},\, \overline{\varrho}_i^{\,2}\bigr]
  \quad \text{[bar\textsuperscript{2}]},
\end{equation}
which makes the steady-state pipe and compressor relations linear or
convex.  The other variables are the supply $x_k$, the gas shed
$y_d \in [0, L_d]$, the pipe flow $f^{\mathrm{g}}_a$ (or the inlet and
outlet flows $f^{\mathrm{in}}_a, f^{\mathrm{out}}_a$ when the pipe
stores gas), the line pack $v^{\mathrm{lp}}_a$, the compressor
throughput $\psi_c$, its fuel $\phi_c$ and the power $p^{\mathrm{el}}_c$
of an electric drive.  The scenario, stage and block indices are
dropped below; every row holds for each $(s,t,b)$ unless stated
otherwise.

\subsection{Gas Balance and Nodal Gas Prices}

Each node balances, in every block, injections against withdrawals:
\begin{multline}
\label{eq:gas_balance}
\sum_{k:\,i_k=i} x_k
+ \sum_{a:\,j_a=i} f^{\mathrm{out}}_a
- \sum_{a:\,i_a=i} f^{\mathrm{in}}_a
+ \sum_{c:\,j_c=i} \psi_c \\
- \sum_{c:\,i_c=i} \chi_c \bigl[(1+\gamma^{\mathrm{ff}}_c)\psi_c
  + \phi_c\bigr]
- \sum_{c:\,i_c=i} (1-\chi_c)\,\psi_c \\
- \sum_{g\in\setG^{\mathrm{g}}:\,i_g=i} h_g\, p_g
+ \sum_{d:\,i_d=i} y_d
= \sum_{d:\,i_d=i} L_d ,
\end{multline}
where $f^{\mathrm{in}}_a = f^{\mathrm{out}}_a = f^{\mathrm{g}}_a$ for a
pipe without line pack, $\chi_c = 1$ for a gas-driven station (its fuel
is drawn at the suction node) and $0$ for an electric drive,
$\gamma^{\mathrm{ff}}_c$ is a fixed fuel fraction (zero unless that fuel
model is selected), $h_g$ is the heat rate of generator~$g$, and $L_d$
the gas load.  Supplies cost $c_k$ and shed gas $c^{\mathrm{fail}}_d$,
both weighted by $\omega_{s,t,b}$ as in~\eqref{eq:opcost}; a supply may
also carry a stage volume
\begin{equation}
\label{eq:gas_volume}
\sum_{b \in \setB_t} \Delta_b\, x_{k,s,t,b} \le V_{k,t},
\end{equation}
which represents an LNG cargo or a pipeline allotment.

The dual $\mu_i$ of~\eqref{eq:gas_balance} is the nodal gas price
[\$/km\textsuperscript{3}], the gas analogue of the LMP.  Because the
fuel of a gas-fired unit enters~\eqref{eq:gas_balance} with coefficient
$-h_g$ on the same column $p_g$ that enters the bus
balance~\eqref{eq:balance}, the optimality conditions of a unit
dispatched strictly between its bounds, with no other active coupling
constraint, give
\begin{equation}
\label{eq:gas_price_link}
\lambda_{n_g} = c_g + h_g\, \mu_{i_g} + \nu^{\mathrm{gate}}_g ,
\end{equation}
where $\lambda_{n_g}$ is the electricity price at its bus and
$\nu^{\mathrm{gate}}_g \geq 0$ the dual of its pressure gate
(Section~\ref{sec:gas_turbine}).  The gas price is therefore set by the
gas supply when the network is uncongested, and by the electricity
price, $(\lambda_{n_g} - c_g)/h_g$, when the delivery to the plant is
limited.  For a flowing pipe without line pack, the reduced-cost
condition of $f^{\mathrm{g}}_a$ gives
$\mu_{j_a} - \mu_{i_a} = \sum_m 2K^e_a t_{a,m}\, \nu^{W}_{a,m}$ plus the
capacity dual, with $\nu^{W}_{a,m} \geq 0$ the duals of the Weymouth
rows~\eqref{eq:gas_weymouth}: the gas price rises along the flow only
when pressure is scarce, i.e.\ when a pressure bound binds somewhere
in the network.

\subsection{Pipes: Weymouth Equation and its Tangent Approximation}

The steady-state flow of a horizontal pipe follows the Weymouth
equation in squared pressure,
$\varpi_{i_a} - \varpi_{j_a} = K_a\, f^{\mathrm{g}}_a |f^{\mathrm{g}}_a|$,
whose constant [bar\textsuperscript{2}/(km\textsuperscript{3}/h)\textsuperscript{2}]
is either given or computed from the SI Weymouth formula of
Menon~\cite[eq.~(2.62)]{menon2005gas}, converted from
(m\textsuperscript{3}/d, kPa) to (km\textsuperscript{3}/h, bar):
\begin{equation}
\label{eq:gas_weymouth_k}
K_a = \frac{G\, T_f\, L_a\, Z}
  {\bigl(3.7435\cdot 10^{-3}\, E_a\, (T_b/P_b)\, D_a^{2.667}\bigr)^2}
  \cdot \frac{24000^2}{10^4},
\end{equation}
with length $L_a$ [km], inner diameter $D_a$ [mm], pipeline efficiency
$E_a$ (default 0.92), gas specific gravity $G$ (0.6), flowing
temperature $T_f$ (288.15~K), compressibility $Z$ (0.9) and base
conditions $T_b = 288.15$~K, $P_b = 101.325$~kPa.

\paragraph{Radial networks.}
When the flow direction of every pipe is known, as in a tree, the flow
is $f^{\mathrm{g}}_a \in [0, \overline{F}_a]$, and the equality is
relaxed to the convex $\varpi_{i_a} - \varpi_{j_a} \geq K_a (f^{\mathrm{g}}_a)^2$.
This relaxation is replaced by $N_a$ tangent planes, the outer
approximation used by Tomasgard et al.~\cite{tomasgard2007gas} and by
Ordoudis et al.~\cite{ordoudis2019gas}, here written with the elevation
term introduced below:
\begin{equation}
\label{eq:gas_weymouth}
\varpi_{i_a} - e^{\zeta_a}\varpi_{j_a}
  \geq K^e_a \bigl(2 t_{a,m}\, \tilde f_a - t_{a,m}^2\bigr),
  \quad m = 1,\dots,N_a .
\end{equation}
Here $\tilde f_a$ is the pipe flow, or the mean
$(f^{\mathrm{in}}_a + f^{\mathrm{out}}_a)/2$ when the pipe stores gas.
With the default uniform layout, $t_{a,m} = (2m-1)\overline{F}_a/(2N_a)$,
the squared-pressure error of the envelope is at most
$K^e_a(\overline{F}_a/2N_a)^2$; an alternative layout places the
tangents uniformly in $\arcsin(\sqrt{K_a} f / \overline{\varrho}_{i_a})$,
which evens out the relative error of the outlet pressure.  The
relaxation allows the LP to drop more pressure than the flow requires.
On a radial network this slack is physically realisable by a pressure
regulator at the pipe end and, since pressure carries no cost, the LP
has no incentive to use it, unless an upper pressure bound downstream
binds.  For every pipe and block, \gtopt{} reports the Weymouth gap
$\varpi_{i_a} - e^{\zeta_a}\varpi_{j_a} - K^e_a \tilde f_a|\tilde f_a|$
and the corresponding outlet-pressure error, and warns when the latter
exceeds 0.1~bar.  The case of Section~\ref{sec:gas_bench} in which this
happens is detected by that report.

\paragraph{Elevation.}
A pipe that climbs from $h_{i_a}$ to $h_{j_a}$ [m] follows the
steady-state equation with gravity~\cite[\S 2.9]{menon2005gas},
\begin{gather}
\varpi_{i_a} - e^{\zeta_a}\varpi_{j_a} = K^e_a\, f|f|,
\label{eq:gas_elevation} \\
\zeta_a = \frac{2\,g\,M\,(h_{j_a} - h_{i_a})}{Z\,R\,T},
\qquad
K^e_a = K_a \frac{e^{\zeta_a} - 1}{\zeta_a},
\nonumber
\end{gather}
with $M = G\, M_{\mathrm{air}}$, $g = 9.80665$~m/s\textsuperscript{2},
$R = 8.314$~J/(mol\,K), and $K^e_a \to K_a$ as $\zeta_a \to 0$.  Both
sides stay linear in $\varpi$, so~\eqref{eq:gas_weymouth} keeps its
form.  The term matters for networks that cross mountain ranges: a
rise of 4000~m gives $\zeta_a \approx 0.63$, i.e.\ the attainable
outlet squared pressure is divided by $e^{\zeta_a} \approx 1.9$ before
friction is accounted for.

\paragraph{Meshed networks and sequential LP.}
In a loop the flow direction is not known in advance and the equality
is non-convex.  A pipe in \emph{linearized} mode therefore has a signed
flow $f^{\mathrm{g}}_a \in [-\overline{F}_a, \overline{F}_a]$ and one
equality per block, linearised at a reference flow $f^0_a$, either as
the secant
$\varpi_{i_a} - e^{\zeta_a}\varpi_{j_a} = K^e_a |f^0_a|\,\tilde f_a$
or as the tangent
$\varpi_{i_a} - e^{\zeta_a}\varpi_{j_a} = 2K^e_a|f^0_a|\,\tilde f_a
- K^e_a f^0_a|f^0_a|$.
Both are exact at $\tilde f_a = f^0_a$.  A sequential LP (SLP) repeats
three steps after the first solve: it updates
$f^0_a \leftarrow f^0_a + \omega (\tilde f_a - f^0_a)$ with damping
$\omega$, re-issues the row coefficients through the same
re-linearisation interface used by the volume-dependent hydro elements
of Section~\ref{sec:hydro} (seepage and discharge limits), and
re-solves.  It stops when the largest
relative flow change falls below a tolerance.  When gas costs tie, the
LP has several optimal vertices between which the iterates can
alternate; a proximal row
$\tilde f_a - d^+_a + d^-_a = f^0_a$, with a small cost on
$d^+_a + d^-_a$, selects among them the vertex nearest $f^0_a$.  At the
fixed point the rows satisfy the exact Weymouth equation; the point is
stationary, but its global optimality is not certified.  The SLP runs
with the monolithic method; under SDDP the rows are re-issued at the
stored reference flows in every forward pass, but the reference flows
are not updated.  A valve is an
open connection with signed flow and $\varpi_{j_c} = \varpi_{i_c}$, and
an arc without pressure coupling, $0 \leq f \leq \overline{F}$, models a
delivery through a parallel system.

\subsection{Line Pack}
\label{sec:gas_linepack}

A pipe with line pack splits its flow into inlet and outlet flows.  The
stored gas $c_a(\varrho_{i_a} + \varrho_{j_a})/2$ is linearised at a
reference pressure $\varrho^{r}_a$ through
$\varrho \approx (\varpi + (\varrho^{r}_a)^2)/(2\varrho^{r}_a)$:
\begin{align}
v^{\mathrm{lp}}_{a,b}
  &= \frac{c_a}{4\varrho^{r}_a}\bigl(\varpi_{i_a,b} + \varpi_{j_a,b}\bigr)
   + \frac{c_a\,\varrho^{r}_a}{2},
   \label{eq:gas_linepack_def} \\
v^{\mathrm{lp}}_{a,b}
  &= v^{\mathrm{lp}}_{a,b-1}
   + \Delta_b\bigl(f^{\mathrm{in}}_{a,b} - f^{\mathrm{out}}_{a,b}\bigr),
   \label{eq:gas_linepack_bal} \\
v^{\mathrm{lp}}_{a,|\setB_t|} &\geq v^{\mathrm{lp},0}_a ,
   \label{eq:gas_linepack_end}
\end{align}
where $c_a = \pi D_a^2 L_a/(4 Z P_b)\,(T_b/T)$ is the geometric
capacity [km\textsuperscript{3}/bar] and $v^{\mathrm{lp},0}_a$ the
initial pack, by default derived from the initial node pressures.  The
pack thus plays the role of a storage whose level is tied to the
pressures, and~\eqref{eq:gas_linepack_bal} has the same form as the
storage balance of Section~\ref{sec:battery}.  The terminal
condition~\eqref{eq:gas_linepack_end} closes a daily cycle; it can be
dropped so that the pack is drawn down to the pressure bounds, e.g.\ in
multi-day supply-interruption studies.  The linearisation overestimates
the pressure by $(\varrho - \varrho^{r}_a)^2/(2\varrho^{r}_a)$, about
5\% of the pack for a $\pm 20$~bar swing.  The line pack is held
within each stage; it is not an SDDP state variable.

\subsection{Compressor Stations}

A station raises the pressure up to a maximum ratio
$\overline{\beta}_c$:
\begin{equation}
\label{eq:gas_ratio}
\varpi_{j_c} \leq \overline{\beta}_c^{\,2}\, \varpi_{i_c},
\qquad
\varpi_{j_c} \geq \varpi_{i_c},
\qquad
0 \leq \psi_c \leq \overline{F}_c .
\end{equation}
The second row is omitted for a station that may be bypassed, which
then acts as a pressure regulator.  Three fuel models are available: no
fuel; a fixed fraction $\phi_c = \gamma^{\mathrm{ff}}_c \psi_c$; and the
model of Saldarriaga~\cite[eqs.~(2.70)--(2.72)]{saldarriaga2018phd},
$\phi_c = \gamma_c(\beta_c^{\alpha} - 1)\psi_c$ with
$\beta_c = \varrho_{j_c}/\varrho_{i_c}$ and
$\alpha = (k-1)/k = 0.2308$ for an isentropic exponent $k = 1.3$.
Written in squared pressure, this fuel is homogeneous of degree one in
$\psi_c$ and zero in $\varpi$, so its first-order Taylor plane at an
operating point $(\psi^0_c, \varpi^0_{i_c}, r^0_c\varpi^0_{i_c})$, with
$r^0_c = (\beta^0_c)^2$, has no constant term:
\begin{multline}
\label{eq:gas_fuel_plane}
\phi_c \geq \gamma_c\bigl((\beta^0_c)^{\alpha} - 1\bigr)\psi_c \\
+ \gamma_c \frac{\alpha}{2}\,\psi^0_c\,(r^0_c)^{\alpha/2 - 1}
  \frac{\varpi_{j_c} - r^0_c\,\varpi_{i_c}}{(\varrho^0_{i_c})^2},
\qquad \phi_c \geq 0 .
\end{multline}
One plane is added per operating ratio $\beta^0_c$, so that several
planes form a piecewise-linear lower envelope.  The inequality is tight
whenever the plane is positive, because fuel costs gas.  When not given,
$\gamma_c$ is derived from the adiabatic compression head
$H = \frac{Z R T_b}{M}\frac{k}{k-1}(\beta^{\alpha}-1)/\eta_c$ and the
fuel $H/(\eta_{\mathrm{drv}}\,\mathrm{LHV})$, with compressor and driver
efficiencies $\eta_c = 0.80$ and $\eta_{\mathrm{drv}} = 0.33$, which
gives $\gamma_c \approx 0.043$.  The fixed fraction defaults to the same
fuel at the maximum ratio, $\gamma_c(\overline{\beta}_c^{\,\alpha}-1)$.
An electric drive does not burn gas; it adds a load
$p^{\mathrm{el}}_c = \kappa^{\mathrm{el}}_c \phi_c \leq
\overline{P}^{\mathrm{el}}_c$ to its bus balance~\eqref{eq:balance},
with $\kappa^{\mathrm{el}}_c =
\mathrm{LHV}\,\eta_{\mathrm{drv}}/\eta_{\mathrm{mot}} \approx 3.3$
MW per km\textsuperscript{3}/h.  Compressor on/off decisions are not
modelled.

\subsection{Gas-Fired Units: Fuel, Pressure Gate and Boosters}
\label{sec:gas_turbine}

A gas-turbine link couples generator $g$ to its delivery node $i_g$.
Its fuel $h_g p_g$ is withdrawn in~\eqref{eq:gas_balance}.  When the
turbine requires a minimum fuel-gas pressure $\underline{\varrho}_g$,
its capability is a concave piecewise-linear function of the squared
pressure:
\begin{align}
p_g &\leq s^{(1)}_g\bigl(\varpi_{i_g}
  - (\underline{\varrho}_g - \delta_g)^2\bigr),
  \label{eq:gas_gate} \\
p_g &\leq o_g\overline{P}_g
  + s^{(2)}_g\bigl(\varpi_{i_g} - \underline{\varrho}_g^{\,2}\bigr),
  \label{eq:gas_derate}
\end{align}
with $s^{(1)}_g = o_g\overline{P}_g/(\underline{\varrho}_g^{\,2}
- (\underline{\varrho}_g - \delta_g)^2)$ and
$s^{(2)}_g = (1 - o_g)\overline{P}_g/((\varrho^{\mathrm{nom}}_g)^2
- \underline{\varrho}_g^{\,2})$.  The gate~\eqref{eq:gas_gate} ramps the
output from zero at $(\underline{\varrho}_g - \delta_g)^2$ to
$o_g\overline{P}_g$ at $\underline{\varrho}_g^{\,2}$, where $\delta_g$
(default 0.5~bar) is a supply tolerance.  The derate
line~\eqref{eq:gas_derate} is added when a nominal pressure
$\varrho^{\mathrm{nom}}_g > \underline{\varrho}_g$ is given, and raises
the output from $o_g\overline{P}_g$ (by default
$o_g = \underline{\varrho}_g/\varrho^{\mathrm{nom}}_g$) to
$\overline{P}_g$; it is dropped if it would make the capability
non-concave.  Because the intercept of~\eqref{eq:gas_gate} is negative,
the gate holds the delivery node above the bottom of the band even when
the unit is off.  Switching the gate off below the band is a
non-convex condition that requires a binary variable; it is not part of
the LP.

A fuel-gas booster is a compressor station between the end of the plant
lateral and the turbine header, typically electrically driven and
bypassable.  It lifts the delivery pressure above the gate at the cost
of motor power, which enters the bus balance, and of compressor fuel,
which enters the gas balance, so the LP decides when boosting is worth
its cost.

\subsection{LNG Supply}

LNG enters the gas network as a supply tagged as regasification, whose
rate is bounded by the send-out capacity and whose stage volume is
bounded by~\eqref{eq:gas_volume}.  The tank inventory itself is the LNG
terminal of Section~\ref{sec:lng}, which is coupled to the gas-fired
units through their heat rates; linking its send-out to a gas node is
not yet supported.  Gas supply described without a network, as priced
fuel bands with offtake budgets, is also common in operation models.
The \texttt{plexos2gtopt} converter keeps such bands either as fuels
with budget constraints or as fuel tanks drained by their generators,
with the leftover valued by a terminal future-cost variable; the tank
is a state variable, so the gas opportunity cost is carried between
SDDP phases like a water value.
